%!TEX root = ../main_thesis.tex

\chapter{Background e stato dell'arte}
\label{cap:background}

Questo capitolo fornisce le basi necessarie a comprendere il resto della tesi.
Si introducono dapprima i concetti fondamentali della programmazione lineare e
della sua estensione intera (Sezioni~\ref{sec:lp-milp}
e~\ref{sec:modelli-ottimizzazione}); si analizzano poi i principali linguaggi e
strumenti di modellazione esistenti (Sezione~\ref{sec:stato-arte}), per
inquadrare il contesto in cui si colloca ROOC; infine si richiamano i concetti
di teoria dei compilatori (Sezione~\ref{sec:teoria-compilatori}) che verranno
applicati nei capitoli successivi, e si delinea il posizionamento del progetto
(Sezione~\ref{sec:posizionamento}).

\section{Programmazione lineare e MILP}
\label{sec:lp-milp}

La \emph{programmazione lineare} (\emph{Linear Programming}, LP) è il ramo
dell'ottimizzazione matematica che studia la minimizzazione (o massimizzazione)
di una funzione obiettivo lineare soggetta a un insieme di vincoli lineari. Un
problema di programmazione lineare può essere espresso nella forma generale

\begin{equation}
\label{eq:lp-generale}
\begin{aligned}
\min_{x} \quad & c^\top x \\
\text{s.t.} \quad & A x \le b, \\
                  & x \ge 0,
\end{aligned}
\end{equation}

dove $x \in \mathbb{R}^n$ è il vettore delle \emph{variabili decisionali},
$c \in \mathbb{R}^n$ il vettore dei coefficienti della funzione obiettivo,
$A \in \mathbb{R}^{m \times n}$ la matrice dei coefficienti dei vincoli e
$b \in \mathbb{R}^m$ il vettore dei termini noti. L'insieme dei punti che
soddisfano tutti i vincoli prende il nome di \emph{regione ammissibile}; quando
non è vuota e la funzione obiettivo è limitata sulla regione, esiste almeno una
\emph{soluzione ottima}, che il problema si propone di determinare~\cite{dantzig1963,bertsimas1997}.

Un risultato classico della teoria afferma che, se esiste una soluzione ottima,
allora ne esiste una in corrispondenza di un \emph{vertice} del poliedro che
descrive la regione ammissibile. Su questa proprietà si fonda il \emph{metodo
del simplesso}, che esplora i vertici della regione spostandosi di volta in
volta verso quello con valore migliore della funzione obiettivo. Il simplesso
opera su una rappresentazione del problema detta \emph{forma standard},

\begin{equation}
\label{eq:lp-standard}
\min_{x} \ c^\top x \quad \text{s.t.} \quad A x = b, \ x \ge 0,
\end{equation}

in cui tutti i vincoli sono uguaglianze e tutte le variabili sono non negative.
La conversione di un problema dalla forma~\eqref{eq:lp-generale} alla
forma~\eqref{eq:lp-standard} è un passaggio puramente meccanico, che verrà
trattato nel Capitolo~\ref{cap:linearizzazione} come una delle fasi della
pipeline di compilazione di ROOC.

In molte applicazioni reali le variabili decisionali non possono assumere valori
frazionari: rappresentano scelte indivisibili (quanti veicoli acquistare) o
decisioni di tipo sì/no (se attivare o meno un impianto). Si parla allora di
\emph{programmazione lineare intera mista} (\emph{Mixed-Integer Linear
Programming}, MILP), in cui si aggiunge il vincolo di integralità su un
sottoinsieme $I \subseteq \{1, \dots, n\}$ delle variabili:

\begin{equation}
\label{eq:milp}
\min_{x} \ c^\top x \quad \text{s.t.} \quad A x \le b, \ x \ge 0, \ x_j \in \mathbb{Z} \ \ \forall j \in I.
\end{equation}

Un caso particolarmente importante è quello delle variabili \emph{binarie}, in
cui $x_j \in \{0, 1\}$, usate per modellare decisioni logiche. L'aggiunta del
vincolo di integralità cambia radicalmente la natura del problema: mentre la
programmazione lineare continua è risolvibile in tempo polinomiale, la
programmazione lineare intera è in generale \emph{NP-hard}~\cite{wolsey1998,nemhauser1988}.
La tecnica più diffusa per risolverla è il \emph{branch-and-bound}, che esplora
in modo sistematico lo spazio delle soluzioni intere risolvendo una successione
di rilassamenti continui; tale tecnica sarà ripresa nel
Capitolo~\ref{cap:solver}, dove si descrive l'estensione del solver
\texttt{microlp} al caso intero e binario.

\section{I modelli di ottimizzazione}
\label{sec:modelli-ottimizzazione}

Tradurre un problema del mondo reale in una delle forme~\eqref{eq:lp-generale}
o~\eqref{eq:milp} significa costruirne un \emph{modello}: identificare le
variabili decisionali, esprimere l'obiettivo come funzione lineare di tali
variabili e codificare i requisiti del problema come vincoli. Questa attività di
\emph{modellazione} è concettualmente distinta dalla \emph{risoluzione}
numerica, ed è proprio il livello a cui opera un linguaggio come ROOC.

Si consideri, a titolo di esempio, il classico problema dello \emph{zaino}
(\emph{knapsack}): dato un insieme di $n$ oggetti, ciascuno con un valore $v_i$
e un peso $w_i$, e uno zaino di capacità $C$, si vuole selezionare il
sottoinsieme di oggetti di valore complessivo massimo il cui peso non superi la
capacità. Introducendo una variabile binaria $x_i \in \{0,1\}$ che indica se
l'oggetto $i$-esimo è selezionato, il modello si scrive

\begin{equation}
\label{eq:knapsack}
\begin{aligned}
\max \quad & \sum_{i=1}^{n} v_i\, x_i \\
\text{s.t.} \quad & \sum_{i=1}^{n} w_i\, x_i \le C, \\
                  & x_i \in \{0, 1\} \quad \forall i \in \{1, \dots, n\}.
\end{aligned}
\end{equation}

Si noti come la struttura del modello sia indipendente dai dati specifici
($n$, i valori $v_i$, i pesi $w_i$, la capacità $C$): la stessa formulazione
descrive un'intera famiglia di istanze. La separazione tra la \emph{struttura}
del modello e i \emph{dati} di una particolare istanza è un principio cardine dei
linguaggi di modellazione, ROOC compreso, e ne motiva i costrutti come gli
iteratori e le sommatorie indicizzate, descritti nel
Capitolo~\ref{cap:linguaggio}. Il problema dello zaino, insieme al problema del
\emph{dominating set} su grafo, farà da esempio guida nel resto della tesi.

\section{Linguaggi e strumenti di modellazione esistenti}
\label{sec:stato-arte}

La modellazione di problemi di ottimizzazione è supportata da un'ampia gamma di
strumenti, che si possono collocare lungo uno spettro che va dalle
rappresentazioni di basso livello, vicine al solver, fino ai linguaggi
algebrici di alto livello, vicini alla notazione matematica.

\paragraph{Formati di scambio di basso livello.}
Formati testuali come il \emph{CPLEX LP format} e il formato \emph{MPS}
descrivono direttamente la matrice dei coefficienti di un modello lineare. Sono
pensati come interfaccia tra strumenti software e solver, non per essere scritti
o letti dall'uomo: non offrono astrazioni quali iteratori, costanti simboliche o
strutture dati, e impongono di esplicitare ogni singolo vincolo. ROOC è in grado
di esportare in formato CPLEX LP, usato come interfaccia verso solver esterni.

\paragraph{Linguaggi algebrici di modellazione.}
I \emph{linguaggi algebrici di modellazione} (\emph{Algebraic Modeling
Languages}, AML) permettono di esprimere i modelli in una forma vicina a quella
matematica, con sommatorie, indici e parametri simbolici. \textbf{AMPL} è il
capostipite di questa famiglia: sviluppato presso i Bell Labs, separa nettamente
modello e dati e si interfaccia con un gran numero di solver~\cite{ampl2003}. Si
tratta però di un prodotto commerciale e a sorgente chiuso. \textbf{GNU
MathProg} (GMPL) ne riprende un sottoinsieme in forma libera, come parte del GNU
Linear Programming Kit~\cite{glpk}. \textbf{GAMS} è un altro storico ambiente
commerciale della stessa categoria.

\paragraph{Linguaggi embedded.}
Una tendenza più recente è quella di offrire la modellazione come libreria
all'interno di un linguaggio di programmazione general-purpose. \textbf{Pyomo}
fornisce un'API a oggetti in Python~\cite{pyomo2017}, mentre \textbf{JuMP},
embedded in Julia, sfrutta la metaprogrammazione del linguaggio ospite per
ottenere al contempo espressività e prestazioni elevate~\cite{jump2017}. Questo
approccio eredita gratuitamente l'ecosistema del linguaggio ospite, ma lega la
sintassi del modello a quella di tale linguaggio.

\paragraph{Linguaggi di modellazione a vincoli.}
\textbf{MiniZinc} è un linguaggio open source per la \emph{programmazione a
vincoli}, indipendente dal solver: un modello MiniZinc viene compilato in un
formato intermedio di più basso livello (FlatZinc) che diversi solver, di tipo
CP, MIP o SAT, possono poi risolvere~\cite{minizinc2007}. Questo schema, in cui
un linguaggio di alto livello viene \emph{compilato} verso una forma intermedia
solver-indipendente, è concettualmente vicino all'architettura di ROOC.

\paragraph{Solver e relative API.}
Va infine distinto il livello del \emph{solver} vero e proprio. Strumenti come
\textbf{Gurobi}~\cite{gurobi} o \textbf{HiGHS}~\cite{highs2018} sono motori di
risoluzione, corredati di API per la costruzione dei modelli, ma non
costituiscono di per sé linguaggi di modellazione. ROOC si appoggia a solver di
questo tipo come back-end intercambiabili (Capitolo~\ref{cap:solver}).

\medskip
Rispetto a questo panorama, ROOC si propone di unire alcune caratteristiche che
raramente coesistono in un unico strumento: una sintassi dichiarativa vicina
alla matematica, ma con un sistema di tipi ricco che include strutture dati come
\emph{grafi}, tuple e iterabili; l'indipendenza dal solver tramite una pipeline
di compilazione esplicita; l'esecuzione interamente lato client tramite
WebAssembly; e un taglio didattico, con la possibilità di ispezionare ogni fase
intermedia della risoluzione. Tali aspetti sono ripresi nella
Sezione~\ref{sec:posizionamento}.

\section{Cenni di teoria dei linguaggi e dei compilatori}
\label{sec:teoria-compilatori}

ROOC è, a tutti gli effetti, un \emph{compilatore}: traduce il linguaggio di
modellazione, di alto livello, in rappresentazioni progressivamente più vicine a
una forma risolvibile da un solver. La pipeline non è vincolata a un'unica forma
di destinazione: l'implementazione attuale arriva fino a un modello lineare in
forma standard, ma la stessa architettura potrebbe in futuro produrre modelli
non lineari, fermarsi a una fase intermedia o aggiungere ulteriori
trasformazioni. È quindi utile richiamare le fasi classiche di un
compilatore~\cite{aho2006}, perché esse forniscono la struttura portante della
pipeline descritta nel resto della tesi.

\paragraph{Analisi lessicale e sintattica.}
Le prime fasi di un compilatore riconoscono la struttura del testo sorgente.
L'\emph{analisi lessicale} suddivide il flusso di caratteri in unità elementari
(\emph{token}); l'\emph{analisi sintattica} (\emph{parsing}) verifica che la
sequenza di token rispetti la grammatica del linguaggio e ne costruisce una
rappresentazione ad albero, l'\emph{albero di sintassi astratta} (\emph{Abstract
Syntax Tree}, AST). La grammatica è tradizionalmente espressa tramite una
grammatica \emph{context-free} (CFG), spesso in notazione BNF.

\paragraph{Parsing Expression Grammar.}
Una formalizzazione alternativa è data dalle \emph{Parsing Expression Grammar}
(PEG)~\cite{ford2004}. A differenza delle CFG, le PEG sono intrinsecamente non
ambigue: la scelta tra alternative è \emph{ordinata}, ossia si prova la prima
alternativa e si passa alla successiva solo in caso di fallimento. Inoltre le
PEG integrano analisi lessicale e sintattica in un'unica grammatica, eliminando
la necessità di un lexer separato. Queste proprietà le rendono particolarmente
adatte a un linguaggio come ROOC; la grammatica concreta, espressa con la
libreria PEST, è descritta nel Capitolo~\ref{cap:parsing}.

\paragraph{Analisi semantica e sistema dei tipi.}
Superata l'analisi sintattica, l'\emph{analisi semantica} verifica che il
programma sia semanticamente corretto, cioè che le regole del linguaggio
vengano rispettate: che le variabili siano dichiarate, che le operazioni siano
applicate a operandi di tipo compatibile, che le funzioni siano invocate con il
numero e il tipo di argomenti corretti. Il \emph{sistema
dei tipi} è lo strumento formale con cui si esprimono e si verificano queste
proprietà~\cite{pierce2002}. In ROOC questa fase comprende anche la
\emph{trasformazione} del modello, in cui costanti, iteratori e sommatorie
vengono risolti ed espansi (Capitolo~\ref{cap:semantica}).

\paragraph{Rappresentazione intermedia e \emph{lowering}.}
I compilatori non traducono direttamente dall'AST al codice finale, ma passano
attraverso una o più \emph{rappresentazioni intermedie} (IR), via via più
vicine alla macchina di destinazione. Il processo di traduzione da una
rappresentazione di livello superiore a una di livello inferiore prende il nome
di \emph{lowering}. Nel caso di ROOC, il \emph{lowering} consiste nella
\emph{linearizzazione} del modello, cioè l'eliminazione dei costrutti non
lineari come \texttt{min}, \texttt{max} e valore assoluto, e nella sua riduzione
alla forma standard, descritte nel Capitolo~\ref{cap:linearizzazione}.

\paragraph{Generazione del codice.}
L'ultima fase di un compilatore tradizionale produce il codice oggetto. In ROOC
il ruolo del codice oggetto è svolto dal modello lineare in forma standard, che
viene infine dato in pasto a un solver per ottenere la soluzione, come discusso
nel Capitolo~\ref{cap:solver}.

\section{Posizionamento di ROOC}
\label{sec:posizionamento}

Alla luce di quanto esposto, ROOC si colloca come linguaggio di modellazione di
alto livello, indipendente dal solver, con alcune scelte progettuali distintive:

\begin{itemize}
  \item una notazione dichiarativa vicina alla matematica, ma arricchita da un
        sistema di tipi che comprende strutture dati di alto livello (come
        grafi, tuple e iterabili) utili a modellare problemi di ottimizzazione
        combinatoria;
  \item un'architettura a compilatore con pipeline esplicita e
        \emph{solver-indipendente}, sulla scia dell'approccio di MiniZinc, ma
        orientata alla programmazione lineare intera;
  \item l'esecuzione interamente lato client tramite la compilazione in
        WebAssembly~\cite{webassembly}, che elimina la necessità di
        un'infrastruttura server e rende lo strumento immediatamente
        accessibile da un browser;
  \item un taglio \emph{didattico}: la possibilità di ispezionare ogni
        rappresentazione intermedia e di visualizzare passo dopo passo
        l'esecuzione del simplesso rende ROOC uno strumento utile non solo per
        risolvere modelli, ma anche per comprenderne il processo di risoluzione.
\end{itemize}

\noindent
I capitoli successivi approfondiscono ciascuno di questi aspetti, a partire
dalla descrizione del linguaggio dal punto di vista dell'utente
(Capitolo~\ref{cap:linguaggio}).
